\documentclass[twocolumn]{article}

% ---- Packages ----
\usepackage[utf8]{inputenc}
\usepackage[T1]{fontenc}
\usepackage{lmodern}
\usepackage{microtype}
\usepackage{graphicx}
\usepackage{hyperref}
\usepackage{xcolor}
\usepackage[margin=2cm]{geometry}
\usepackage{enumitem}

\hypersetup{
  colorlinks=true,
  linkcolor=blue!60!black,
  citecolor=blue!60!black,
  urlcolor=blue!60!black
}

% ---- Title ----
\title{BIME: A Browser-Based Molecular Editor with Built-in Substructure Intelligence}

\author{
  Syed Asad Rahman\thanks{Corresponding author: \texttt{s9asad@gmail.com}} \\
  BioInception PVT LTD, Cambridge, UK
}

\date{}

\begin{document}

\maketitle

% ====================================================================
\begin{abstract}
BIME (BioInception Molecular Editor) is a browser-based chemical
structure editor written in pure JavaScript with zero runtime
dependencies.  It embeds the SMSD Pro engine for maximum common
substructure (MCS) detection, substructure search, and molecular
fingerprint computation entirely in the client, eliminating
server round-trips.  BIME also implements IUPAC 2013 CIP Rules~1--4a
for R/S and E/Z stereodescriptor assignment.  Its single-file
architecture enables fully offline deployment in educational and
resource-constrained settings.  BIME joins a rich ecosystem of
community tools and is freely available under the Apache
License~2.0 at \url{https://asad.github.io/bime-dist/}.
\end{abstract}

\noindent\textbf{Keywords:} molecular editor, maximum common
substructure, molecular fingerprints, Cahn--Ingold--Prelog,
JavaScript, cheminformatics

% ====================================================================
\section{Introduction}

Browser-based molecular editors are essential components of modern
cheminformatics applications.  The community benefits from several
excellent tools---JSME~\cite{Bienfait2013}, Ketcher~\cite{Ketcher2024},
and ChemDoodle Web Components~\cite{Burger2015}---each serving distinct
needs.  Server-side toolkits including RDKit~\cite{Landrum2006},
the Chemistry Development Kit (CDK)~\cite{Steinbeck2003}, and
Open Babel~\cite{OBoyle2011} provide the algorithmic backbone for
many cheminformatics workflows.  WebAssembly ports such as
RDKit.js~\cite{RDKitJS2024} have recently brought some of this
capability to the browser, though at the cost of multi-megabyte
downloads and asynchronous loading.

A gap remains for applications requiring a lightweight, zero-dependency
editor with built-in substructure intelligence that can run entirely
offline---from a \texttt{file://} URI on a USB drive, for example.
BIME addresses this gap.  Rather than replacing existing tools, BIME
contributes a complementary approach: embedding the SMSD Pro engine
directly into the editor so that MCS, substructure search, fingerprints,
and stereochemistry assignment all run in a single JavaScript file
with no server, no WebAssembly module, and no network connection
required.

BIME builds on work from BioInception, which also develops
SMSD Pro, ReactionDecoder~\cite{Rahman2016}, and
EC-BLAST~\cite{Rahman2014} for reaction analysis and enzyme
classification.

% ====================================================================
\section{Implementation}

\subsection{Architecture}

BIME consists of approximately 25\,000 lines of ES5 JavaScript
organised into 23 modules following the Immediately Invoked Function
Expression (IIFE) pattern, concatenated into a single distributable
file.  No transpiler, bundler, or package manager is required.
The editor runs in all modern browsers and from \texttt{file://}
URIs without a web server.

The core data structure represents molecules as atom--bond graphs
with 2D coordinates.  Atoms carry element symbol, charge, isotope,
chirality, and hydrogen count; bonds carry type and stereo markers.
SMILES~\cite{Weininger1988} (including reaction SMILES), SMARTS,
MOL V2000, and MOL V3000 are supported for import and export.

\subsection{SMSD Pro Engine}

The SMSD Pro engine, embedded in BIME, provides three capabilities:
(1)~substructure search via VF2++~\cite{Juttner2018} with
Neighbourhood Label Frequency pruning;
(2)~maximum common substructure detection via a multi-level funnel
including McSplit partition refinement~\cite{McCreesh2017} and
Bron--Kerbosch~\cite{Bron1973} with Tomita pivoting; and
(3)~molecular fingerprints---ECFP/FCFP~\cite{Rogers2010},
path-based, and topological torsion---with six similarity metrics
(Tanimoto, Dice, cosine, Soergel, count-Tanimoto, count-Dice).

SMSD Pro is an independent JavaScript implementation that builds on
the algorithmic ideas first described in the original Small Molecule
Subgraph Detector (SMSD)~\cite{Rahman2009}.  It is not a fork of
the original Java codebase but a ground-up reimplementation with
modern algorithms.

\subsection{CIP Stereochemistry}

The CIP module implements Rules~1--4a of the IUPAC 2013
recommendations~\cite{Favre2013} using digraph expansion with
phantom-node duplication for double bonds and ring back-edges.
R/S labels are assigned to tetrahedral centres and E/Z labels to
restricted double bonds.

\subsection{2D Layout}

The layout engine follows the Structure Diagram Generation
architecture~\cite{Helson1999}: ring perception via Smallest Set
of Smallest Rings (SSSR), ring-system fusion, zigzag chain
extension, and overlap resolution.  Forty-five scaffold templates
provide instant coordinates for common ring systems.  A two-phase
simulated-annealing crossing reduction with SMACOF stress
majorisation~\cite{deLeeuw1977} improves diagram clarity for
polycyclic structures.

% ====================================================================
\section{Key Features}

Five capabilities distinguish BIME from existing browser-based
editors (see Figure~\ref{fig:workbench}):

\begin{enumerate}[nosep]
\item \textbf{Client-side MCS and similarity.}  Pairwise MCS
  computation completes in a median of 38\,ms for typical drug-sized
  molecules (20--50 heavy atoms), with results highlighted
  directly on the canvas.
\item \textbf{Client-side fingerprints.}  ECFP, FCFP, path-based,
  and topological torsion fingerprints with six similarity coefficients,
  all computed in the browser.
\item \textbf{CIP stereodescriptor assignment.}  R/S and E/Z labels
  via IUPAC 2013 Rules~1--4a with digraph expansion, displayed
  interactively as the user draws.
\item \textbf{Zero-dependency offline deployment.}  A single HTML
  file runs from \texttt{file://} URIs with no server, no WASM
  module, and no network.
\item \textbf{Structure library.}  552 curated structures (drugs,
  amino acids, nucleotides, reagents) for rapid access.
\end{enumerate}

\begin{figure}[t]
\centering
\fbox{\parbox{0.9\columnwidth}{\centering\small
  [Screenshot of BIME workbench showing the editor canvas,\\
  toolbar, and SMSD Pro panel with MCS/fingerprint controls]\\[4pt]
  Interactive demo:
  \url{https://asad.github.io/bime-dist/workbench.html}}}
\caption{The BIME workbench interface.  The left panel provides
  drawing tools and ring templates; the right panel exposes SMSD Pro
  functionality including MCS computation, substructure search, and
  fingerprint similarity scoring.}
\label{fig:workbench}
\end{figure}

% ====================================================================
\section{Limitations and Future Work}

Several limitations should be noted honestly.  The CIP implementation
covers Rules~1--4a; Rules~4b--5 (pseudoasymmetric centres) are not
yet supported.  The crossing-reduction heuristic does not fully
eliminate crossings for all bridged polycyclic systems.  For molecules
exceeding approximately 200 heavy atoms, server-side toolkits such as
RDKit or CDK will offer superior performance.  The distributed
JavaScript is obfuscated; readable source is available from the author
on request for academic verification.  Systematic benchmarking against
standard datasets (GDB, COCONUT) is planned.  Future work will extend
CIP coverage to Rules~4b--5 and explore WebAssembly acceleration for
large molecules.

% ====================================================================
\section*{Availability}

BIME is freely available under the Apache License~2.0.
Live demo: \url{https://asad.github.io/bime-dist/workbench.html}.
Distribution: \url{https://github.com/asad/bime-dist}.
Readable source is available from the author on request.

% ====================================================================
\section*{Acknowledgements}

The author acknowledges the algorithmic foundations laid by the original
SMSD toolkit~\cite{Rahman2009} and the SMSD community (2008--2018).
The author is grateful to the broader cheminformatics community,
including the developers of JSME, Ketcher, RDKit, CDK, and Open Babel,
whose tools and published algorithms have advanced the field and
informed this work.

% ====================================================================
\begin{thebibliography}{99}

\bibitem{Rahman2009}
Rahman, S.A., Bashton, M., Holliday, G.L., Schrader, R., Thornton, J.M.
\newblock Small Molecule Subgraph Detector (SMSD) toolkit.
\newblock \emph{J. Cheminform.}, 1:12, 2009.
\newblock \href{https://doi.org/10.1186/1758-2946-1-12}{doi:10.1186/1758-2946-1-12}

\bibitem{Rahman2014}
Rahman, S.A., Cuesta, S.M., Furnham, N., Holliday, G.L., Thornton, J.M.
\newblock EC-BLAST: a tool to automatically search and compare enzyme reactions.
\newblock \emph{Nat. Methods}, 11(2):171--174, 2014.
\newblock \href{https://doi.org/10.1038/nmeth.2803}{doi:10.1038/nmeth.2803}

\bibitem{Rahman2016}
Rahman, S.A., Torrance, G., Baldacci, L., {et al.}
\newblock Reaction Decoder Tool (RDT): extracting features from chemical reactions.
\newblock \emph{Bioinformatics}, 32(13):2065--2066, 2016.
\newblock \href{https://doi.org/10.1093/bioinformatics/btw096}{doi:10.1093/bioinformatics/btw096}

\bibitem{Bienfait2013}
Bienfait, B., Ertl, P.
\newblock JSME: a free molecule editor in JavaScript.
\newblock \emph{J. Cheminform.}, 5:24, 2013.
\newblock \href{https://doi.org/10.1186/1758-2946-5-24}{doi:10.1186/1758-2946-5-24}

\bibitem{Ertl2010}
Ertl, P.
\newblock Molecular structure input on the web.
\newblock \emph{J. Cheminform.}, 2:1, 2010.
\newblock \href{https://doi.org/10.1186/1758-2946-2-1}{doi:10.1186/1758-2946-2-1}

\bibitem{Burger2015}
Burger, M.C.
\newblock ChemDoodle Web Components: HTML5 toolkit for chemical graphics,
  interfaces, and informatics.
\newblock \emph{J. Cheminform.}, 7:35, 2015.
\newblock \href{https://doi.org/10.1186/s13321-015-0085-3}{doi:10.1186/s13321-015-0085-3}

\bibitem{Ketcher2024}
EPAM Systems.
\newblock Ketcher: web-based molecular editor.
\newblock \url{https://github.com/epam/ketcher}, 2024.

\bibitem{Steinbeck2003}
Steinbeck, C., Han, Y., Kuhn, S., {et al.}
\newblock The {Chemistry Development Kit (CDK)}: an open-source {Java} library
  for chemo- and bioinformatics.
\newblock \emph{J. Chem. Inf. Comput. Sci.}, 43(2):493--500, 2003.
\newblock \href{https://doi.org/10.1021/ci025584y}{doi:10.1021/ci025584y}

\bibitem{Landrum2006}
Landrum, G.
\newblock {RDKit}: open-source cheminformatics.
\newblock \url{https://www.rdkit.org/}, 2006.

\bibitem{OBoyle2011}
O'Boyle, N.M., Banck, M., James, C.A., {et al.}
\newblock Open Babel: an open chemical toolbox.
\newblock \emph{J. Cheminform.}, 3:33, 2011.
\newblock \href{https://doi.org/10.1186/1758-2946-3-33}{doi:10.1186/1758-2946-3-33}

\bibitem{RDKitJS2024}
RDKit.js Contributors.
\newblock {RDKit.js}: {RDKit} for {JavaScript} via {WebAssembly}.
\newblock \url{https://github.com/rdkit/rdkit-js}, 2024.

\bibitem{McCreesh2017}
McCreesh, C., Prosser, P.
\newblock A partitioning algorithm for maximum common subgraph problems.
\newblock In \emph{Proc. IJCAI}, pages 712--719, 2017.
\newblock \href{https://doi.org/10.24963/ijcai.2017/99}{doi:10.24963/ijcai.2017/99}

\bibitem{Juttner2018}
J\"uttner, A., Madarasi, P.
\newblock {VF2++}---An improved subgraph isomorphism algorithm.
\newblock \emph{Discrete Appl. Math.}, 242:69--81, 2018.
\newblock \href{https://doi.org/10.1016/j.dam.2018.02.018}{doi:10.1016/j.dam.2018.02.018}

\bibitem{Bron1973}
Bron, C., Kerbosch, J.
\newblock Algorithm 457: finding all cliques of an undirected graph.
\newblock \emph{Commun. ACM}, 16(9):575--577, 1973.
\newblock \href{https://doi.org/10.1145/362342.362367}{doi:10.1145/362342.362367}

\bibitem{Rogers2010}
Rogers, D., Hahn, M.
\newblock Extended-connectivity fingerprints.
\newblock \emph{J. Chem. Inf. Model.}, 50(5):742--754, 2010.
\newblock \href{https://doi.org/10.1021/ci100050t}{doi:10.1021/ci100050t}

\bibitem{Weininger1988}
Weininger, D.
\newblock {SMILES}, a chemical language and information system.
\newblock \emph{J. Chem. Inf. Comput. Sci.}, 28(1):31--36, 1988.
\newblock \href{https://doi.org/10.1021/ci00057a005}{doi:10.1021/ci00057a005}

\bibitem{Favre2013}
Favre, H.A., Powell, W.H.
\newblock \emph{Nomenclature of Organic Chemistry: IUPAC Recommendations and
  Preferred Names 2013}.
\newblock Royal Society of Chemistry, 2013.
\newblock \href{https://doi.org/10.1039/9781849733069}{doi:10.1039/9781849733069}

\bibitem{Helson1999}
Helson, H.E.
\newblock Structure diagram generation.
\newblock \emph{Rev. Comput. Chem.}, 13:313--398, 1999.
\newblock \href{https://doi.org/10.1002/9780470125908.ch6}{doi:10.1002/9780470125908.ch6}

\bibitem{deLeeuw1977}
de~Leeuw, J.
\newblock Applications of convex analysis to multidimensional scaling.
\newblock In \emph{Recent Developments in Statistics}, pages 133--145.
  North-Holland, 1977.

\bibitem{Morgan1965}
Morgan, H.L.
\newblock The generation of a unique machine description for chemical structures.
\newblock \emph{J. Chem. Doc.}, 5(2):107--113, 1965.
\newblock \href{https://doi.org/10.1021/c160017a018}{doi:10.1021/c160017a018}

\end{thebibliography}

\end{document}
